\section{Запрос на получение сертификата}\label{FMT.CSR}

\subsection{Структура}\label{FMT.CSR.Structure}

Запрос на получение сертификата описывается типом 
\code{CertificationRequest}, который определен в СТБ~34.101.17. 
%
Компонент~\code{subject} запроса должен быть составлен в соответствии
с требованиями раздела~\ref{ENTITIES.Name}.

В запрос могут быть включены атрибуты, которые описываются 
типом~\code{Attribute}, также определенным в СТБ~34.101.17. 
Атрибут представляет собой пару <<идентификатор типа 
(компонент~\code{type})~--- значение (компонент~\code{value})>>.

Разрешается использовать следующие атрибуты:
\begin{enumerate}
\item
Атрибут~\code{challengePassword}. Определен в~\cite{PKCS9}.
%
Идентификатор типа равняется \verb|{1 2 840 113549 1 9 7}|.
Значение~--- строка типа \verb|UTF8String(SIZE(1..255))|.

\item
Атрибут~\code{extensionRequest}. Определен в~\cite{PKCS9}.
%
Идентификатор типа равняется~\verb|{1 2 840 113549 1 9 14}|.
Значение~--- структура типа~\verb|Extensions|, определенного в СТБ~34.101.19.
В компонентах \verb|Extensions| указываются расширения сертификата,
которые планируется перенести в сертификат.

\item
Атрибут \code{certificateValidity}.
Идентификатор типа равняется \verb|bpki-at-certificateValidity|
(определен в приложении~\ref{ASN1}). Значение~--- структура 
типа~\verb|Validity|, определенного в СТБ~34.101.19. 
%
В компонентах \code{notBefore} и \code{notAfter} контейнера~\verb|Validity| 
указываются даты начала и окончания действия сертификата,
рекомендуемые для переноса в сертификат.
%
УЦ может принять рекомендации полностью, частично или вообще проигнорировать. 
\end{enumerate}

В запросах к ПУЦ могут содержаться дополнительные атрибуты.

\subsection{Атрибут~\code{challengePassword}}\label{FMT.CSR.CP}

Строка~--- значение атрибута~\code{challengePassword} состоит из двух частей:
\begin{enumerate} 
\item
Билет выпуска сертификата. Используется в сценарии~\code{Enroll3} 
процесса~\code{Enroll} (см.~\ref{PROCESSES.Enroll}), 
доказывая полномочия на выпуск.

\item
Информационная строка, которую требуется передать УЦ.
Например, реквизиты платежного документа об оплате услуги (процесса).
\end{enumerate}

Первая часть атрибута представляет собой строку длины~$32$, $48$ или~$64$ 
в алфавите  $\{\str{0},\str{1},\ldots,\str{F}\}$ (см.~\ref{CRYPTO.Pwd}).
Длина второй части не должна превышать~$128$. 
Любая из частей может быть опущена. 
Порядок частей не контролируется.
Части одного типа не должны повторяться.

Билету выпуска должен предшествовать префикс~\str{/EPWD:},
информационной строке~--- префикс~\str{/INFO:},
например: \str{/EPWD:01234...EF/INFO:SN112358}.

\subsection{Атрибут~\code{extensionRequest}}\label{FMT.CSR.ER}

В атрибуте \code{extensionRequest} могут быть указаны следующие 
расширения сертификата:
\begin{enumerate}
\item 
\code{ExtKeyUsage}. Расширение указывается в тех случаях, когда будущий субъект
планирует выступать в роли сервера/клиента терминального режима. Соответственно,
расширение может содержать только
идентификаторы~\verb|bpki-eku-serverTM|/\verb|bpki-eku-clientTM|
(см.~\ref{FMT.Ext.EKU}).

\item 
\code{SubjectAltName}. В расширении указываются дополнительные идентификационные
атрибуты (см.~\ref{ENTITIES.SAN}). Расширение должны включать в свои запросы
TLS-сервер и КА, могут включать другие стороны.

\item
\code{CertificatePolicies}. В расширении указываются идентификаторы ролей
будущего субъекта сертификата (см.~\ref{FMT.Ext.CP}). Расширение должны включать
в свои запросы конечные участники и не должны включать ИУЦ и ПУЦ.
\end{enumerate}

В запросах к ПУЦ атрибут~\code{extensionRequest} может содержать дополнительные
расширения.
